\documentclass{article}
\usepackage{geometry}
\usepackage{graphicx} % Required for inserting images
\usepackage{amsmath, amsthm, amssymb}
\usepackage{parskip}
\newgeometry{vmargin={15mm}, hmargin={24mm,34mm}}
\theoremstyle{definition} 
\newtheorem{definition}{Definition}

\newtheorem{theorem}{Theorem}[section]
\newtheorem{lemma}[theorem]{Lemma}
\newtheorem{proposition}[theorem]{Proposition}
\newtheorem{example}[theorem]{Example}
\newtheorem{corollary}{Corollary}[theorem]
\newcommand{\nil}{\operatorname{nil}}
\newcommand{\ann}{\operatorname{Ann}}
\newcommand{\ass}{\operatorname{Ass}}

\title{Schemes}
\date{April 2025}
\author{Boran Erol}

\begin{document}

\maketitle

\section{Pushforward}

\begin{lemma}
    The direct image of a (pre)sheaf on $ X $ is a (pre)sheaf on $ Y $.
\end{lemma}
\begin{proof}
    Let $ \pi: X \to Y $ be a continuous map.

    Let $ \mathcal{F} $ be a sheaf on $ X $.
    
    Let's show that the pushforward sheaf on $ Y $ satisfies
    the sheaf axioms.
    
    Let $ V_{1} \subseteq V_{2} $ in $ Y $.
    Then, $ \pi^{-1}(V_{1}) \subseteq V_{2} $ in $ X $.
    Thus, we have an induced restriction map using the restriction map from $ X $.
    
    Similarly, we have that the restriction triangle is commutative using the commutativity on $ X $.
    
    Thus, $ \pi_{*} \mathcal{F} $ is a presheaf.
    
    Now, let $ V $ be an open subset of $ Y $ and let $ \{ V_{i} \}_{i \in I } $
    be an open cover of $ V $. Notice that $ \{ \pi^{-1}(V_{i}) \}_{i \in I } $ is an
    open cover of $ \pi^{-1}(V) $.
    
    Let $ f_{1}, f_{2} \in \pi_{*} \mathcal{F}(V) = \mathcal{F}(\pi^{-1}(V))$ .
    But then this is simply the gluability axiom and identity axioms on $ X $.
\end{proof}


\begin{lemma}
    Let $ \pi: X \to Y $ be a continuous map and $ \mathcal{F} $ be a sheaf of sets,
    rings or $ A $-modules on $ X $. Let $ p \in X $ and $ q \in Y $ such that $ \pi(x) = y $.
    Then, the pushforward sheaf induces a map on the stalks
\end{lemma}
\begin{proof}
    Let $ (g, V) $ be a germ of $ (\pi_{*} \mathcal{F})_{q} $.

    Let $ U = \pi^{-1}(V) $.

    Then, $ q \in V $ and $ g \in \pi_{*} \mathcal{F}(V) = \mathcal{F}(U)$.

    Map $ (g, V) $ to $ (g, U) $.

    We need to show that this assignment is well-defined.


\end{proof}

Let's now prove the stame statement using the colimit definition of stalks.

\begin{proof}
    
\end{proof}



\end{document}
